\section{Concluding remarks and open questions}

The results above show that one real parameter, determined by the cusp
density, governs the large-genus behavior of $V_{g,n}$ in every regime.
As $n/g$ runs from zero to infinity, the saddle $X(Y_{n/g})$ moves from
the fixed-cusp regime of Mirzakhani and Zograf to the Bessel singularity
$x_c$ of the fixed-genus theory of Manin and Zograf. The volume ratios,
the boundary factor $\mathfrak b_Y$, the short-curve intensities $w_k$
and the tight-length scale $\alpha$ interpolate in the same way.

The coefficients of the normalized volume expansion are analytic in
$Y$ on a neighborhood of the compact interval $[0,y_0]$. This is the
precise sense in which the two endpoint regimes fit into one expansion.
It avoids the loss of uniformity of the fixed-$n$ expansion at
$n\asymp\sqrt g$, while retaining an error bound independent of $n$.

The short geometry nevertheless changes at the scale $\sqrt g$.
For $n=o(\sqrt g)$, Shen and Wu proved a uniform spectral
gap~\cite[Theorem~3]{SW}. At $n/\sqrt g\to c\in(0,\infty)$, the
primitive length spectrum is a Poisson background together with
independent clusters formed by two-cusp pairs of pants
(Theorem~\ref{thm:critical-cluster}). Above this scale, short curves
bounding genus-zero subsurfaces containing cusps have the statistics of
Theorem~\ref{thm:pants-poisson}, and the spectral gap closes
(Corollary~\ref{cor:spectral}).

The proof rests on three features that may be useful beyond
Weil--Petersson volumes. Along the string flow, the alternating
$\kappa_1$--$\psi$ formula becomes a recursion with nonnegative
coefficients, so that comparison and monotonicity arguments hold
uniformly in the deformation parameter. The deformation ends at pure
$\psi$ intersection numbers, where Painlev\'e~I fixes the integration
constants without any resurgence assumption. Finally, the positive
convolution factor of Lemma~\ref{lem:convolution} turns estimates for
generating functions into a local limit theorem that holds uniformly
from bounded cusp counts up to the critical cusp weight. The geometric
consequences combine these volume estimates with the integration
formulas and short-curve arguments credited in
Section~\ref{sec:geometry}.

\subsection{Cusp weights and the JT spectral edge}\label{subsec:jt-edge}

Saad, Shenker and Stanford related the genus expansion of
Jackiw--Teitelboim (JT) gravity to a matrix integral, with
Weil--Petersson volumes as its geometric input~\cite{SSS}.
The parameter $Y$ describes a spectral edge in its defect deformation,
introduced by Maxfield--Turiaci and by Witten~\cite{MT,WittenDef}.
We record the normalization to relate the
uniform estimates to that model. A defect of opening angle $2\pi\alpha$
is a cusp at $\alpha=0$~\cite[Section~2]{ET}.
Let $\lambda_{\mathrm{def}}$ be its weight, distinct from the cusp
density $\lambda=n/g$, and use the convention $2\gamma=1$ of
Maxfield--Turiaci~\cite{MT}. Their string equation at $\alpha=0$
and its branch connected to the undeformed edge are
\[
 \frac{\sqrt{E_0}}{2\pi}I_1(2\pi\sqrt{E_0})
       +\lambda_{\mathrm{def}}=0,\qquad
 E_0=-\frac{Y}{\pi^2},\qquad x=2\pi^2\lambda_{\mathrm{def}}.
\]
The Bessel expressions are interpreted by their entire power series.
Substitution gives exactly $X(Y)=x$. Moreover, the definition
\eqref{eq:H8} gives, for $g\ge2$,
\[
 \sum_{n\ge0}V_{g,n}\frac{\lambda_{\mathrm{def}}^n}{n!}
       =(2\pi^2)^{3g-3}\Phi_g(x).
\]
Thus a fixed cusp count is a coefficient of the defect-weighted
amplitude, whereas $Y$ parametrizes its spectral edge.

The genus-zero density, with its overall topological factor removed,
can be written using our kernel. Specializing the Abel integral in
\cite[equation~(3.50)]{MT} to $\alpha=0$ gives
\[
 \rho_0(E)=\frac1{4\pi}\int_{E_0}^{E}
        \frac{I_0(2\pi\sqrt u)}{\sqrt{E-u}}\,\mathrm du
   =\frac{\sqrt q}{2\pi}K_Y(-2\pi^2q),\qquad q=E-E_0\ge0.
\]
For the second equality set $u=E_0+q(1-s^2)$ and use
\eqref{eq:P1}. At $Y=0$ this is the JT density
$\sinh(2\pi\sqrt E)/(4\pi^2)$. At $Y=y_0$, put
$b=-X''(y_0)>0$. Since $a(y_0)=0$, expansion of $K_{y_0}$ gives
\[
 \rho_0(E_0+q)=\frac{\pi b}{3}q^{3/2}+O(q^{5/2}).
\]
The square-root edge has become a cubic critical edge, the local
geometry associated with the $(2,3)$ critical point. This identifies
the physical meaning of the same endpoint at which the normalized
intersection recursion reduces to pure $\psi$ intersections.

The exponent in the large-genus action agrees with this critical
geometry. With $\delta=y_0-Y$, the endpoint
expansions of Section~\ref{the-stationary-parameter} give
\[
 x_c-x=\frac b2\delta^2(1+O(\delta)),\qquad
 \mathcal Z=\frac{2\sqrt3\,b}{5}\delta^{5/2}(1+O(\delta))
       =\frac{2^{9/4}\sqrt3}{5b^{1/4}}
          (x_c-x)^{5/4}\bigl(1+O(\sqrt{x_c-x})\bigr).
\]
The marked expansion in Theorem~\ref{thm:full-expansion} uses the
full action $\mathcal Z(Y)$ and remains uniform as $Y$ approaches
$y_0$. Replacing the action by its leading critical power would instead
introduce a multiplicative factor $\exp(O(g\sqrt{x_c-x}))$ for fixed
marking count. That replacement is not uniform in an arbitrary joint
limit. In formal genus counting, the corresponding balance is between
the string coupling $\hbar$ and $(x_c-x)^{5/4}$; the genus index $g$
is not the inverse string coupling. A double-scaled sum over all genera
requires a separate summation prescription and estimates beyond the
large-genus coefficient asymptotics proved here.


\subsection{Open questions}
We conclude with some questions that the present methods leave open.

\emph{Actual Gevrey-1 remainders.} The most direct analytic strengthening would replace the squared factorial in Theorem~\ref{thm:intro-absolute} by a single factorial, with constants uniform in the cusp count and truncation order $K\le cg$. A sufficiently small linear truncation would then give error $O(\mathrm e^{-ag})$. The pure-$\psi$ endpoint already has this order of remainder (Subsection~\ref{sec:endpoint-expansion}), but transporting it across the whole parameter interval requires sharper estimates for the actual recursion. In the present proof the squared factorial comes from the degree paths in the profile recursion (Subsection~\ref{subsec:all-orders}); bounds on the formal coefficients alone would not settle the question, since the actual remainders must be controlled as well.

\emph{Rate for adjacent ratios.} Theorem~\ref{thm:coefficient} has the
uniform error $O(g^{-1/8})$, and $O_\varepsilon(g^{-1/4})$ for
$k\ge\varepsilon g$; neither exponent is claimed to be optimal
(Appendix~\ref{sec:quantitative-coefficients}). What is the optimal
uniform rate?

\emph{Other backgrounds.} The derivation of the recursion
\eqref{eq:H9} uses only the string, dilaton and
Dijkgraaf--Verlinde--Verlinde equations, and applies to other
backgrounds of Kontsevich--Witten times; the higher Weil--Petersson
volumes of Kaufmann, Manin and Zagier arise in this way~\cite{KMZ,ManZ}.
For which backgrounds does the deformed recursion have nonnegative
coefficients, and does the uniform analysis of this paper then extend?

\emph{Several boundary components.} Can one obtain relative asymptotics
for mixed intersection numbers when both the number of nonzero degrees
and their sum grow with the genus, uniformly over their allocation among
the markings? Proposition~\ref{prop:bounded-boundaries} treats a fixed
number of bounded nonzero lengths; it does not provide this control of
individual mixed intersections.

\emph{Geometry and spectrum.} Theorem~\ref{thm:tight-spectrum} determines
the tight spectrum at every cusp density, while
Theorem~\ref{thm:critical-cluster} identifies the pants clusters in the
ordinary spectrum when $n/\sqrt g$ has a positive finite limit.
What is the joint distribution
of ordinary geodesics in the same homotopy class after filling the
cusps, when $n$ is proportional to $g$? Our spectral results are upper bounds
for small eigenvalues; their limiting distribution would require
additional geometric input.
